\section{USAMO 2004}
        \begin{enumerate}
        	\item (\textit{USAMO 2004}) Let $ABCD$ be a \href{https://artofproblemsolving.com/wiki/index.php?title=Quadrilateral}{quadrilateral} circumscribed about a circle, whose interior and exterior angles are at least $60 ^{\circ}$.  Prove that


 
$\frac{1}{3} | AB^3 - AD^3 | \le | BC^3 - CD^3 | \le 3 |AB^3 - AD^3 |$


 


 When does equality hold?


 



\textbf{Solution 1}

It suffices to show that 
 \[|BC^3-CD^3| \leq 3|AB^3-AD^3|,\] because the other side of the inequality follows by symmetry. By Pitot's Theorem, we have $AB-AD=BC-CD$. WLOG, let $AB \geq AD$. Then we have that $BC \geq CD$, so we must show $BC^3-CD^3 \leq 3(AB^3-AD^3)$. If $AB=AD$, we are done. Therefore, we will assume for the rest of the proof that $AB>AD$, $BC>CD$.

 Then we would like to show that 
 \[BC^2+BC\cdot CD+CD^2 \leq 3(AB^2+AB\cdot AD+AD^2).\]By the Law of Cosines, we have that 
 \[BD^2=BC^2+k_1(BC\cdot CD)+CD^2=AB^2+k_2(AB\cdot AD)+AD^2,\] where $-1 \leq k_1, k_2 \leq 1$ (this is because $60^{\circ} \leq \angle BAD,\angle BCD \leq 120^{\circ}$). Therefore, 
 \[3(BC^2-BC\cdot CD+CD^2) \leq 3(BC^2+k_1(BC\cdot CD)+CD^2)=3(AB^2+k_2(AB\cdot AD)+AD^2)\leq 3(AB^2+AB\cdot AD+AD^2).\] Then we wish to show that 
 \[3(BC^2-BC\cdot CD+CD^2)-(BC^2+BC\cdot CD+CD^2) \geq 0 \rightarrow 2(BC-CD)^2 \geq 0,\] and we are done. $\blacksquare$

 



\textbf{Solution 2}

By a well-known property of tangential quadrilaterals, the sum of the two pairs of opposite sides are equal; hence $a + c = b + d \Rightarrow a - b = d - c \Rightarrow |a - b| = |d - c|$ Now we factor the desired expression into $\frac {|d - c|(c^2 + d^2 + cd)}{3} \le|a - b|(a^2 + b^2 + ab)\le 3|d - c|(c^2 + d^2 + cd)$. Temporarily discarding the case where $a = b$ and $c = d$, we can divide through by the $|a - b| = |d - c|$ to get the simplified expression $(c^2 + d^2 + cd)/3\le a^2 + b^2 + ab\le 3(c^2 + d^2 + cd)$.

 Now, draw diagonal $BD$. By the law of cosines, $c^2 + d^2 + 2cd\cos A = BD$. Since each of the interior and exterior angles of the quadrilateral is at least 60 degrees, we have that $A\in [60^{\circ},120^{\circ}]$. Cosine is monotonically decreasing on this interval, so by setting $A$ at the extreme values, we see that $c^2 + d^2 - cd\le BD^2 \le c^2 + d^2 + cd$. Applying the law of cosines analogously to $a$ and $b$, we see that $a^2 + b^2 - ab\le BD^2 \le a^2 + b^2 + ab$; we hence have $c^2 + d^2 - cd\le BD^2 \le a^2 + b^2 + ab$ and $a^2 + b^2 - ab\le BD^2 \le c^2 + d^2 + cd$. 

 We wrap up first by considering the second inequality. Because $c^2 + d^2 - cd\le BD^2 \le a^2 + b^2 + ab$, $\text{RHS}\ge 3(a^2 + b^2 - ab)$. This latter expression is of course greater than or equal to $a^2 + b^2 + ab$ because the inequality can be rearranged to $2(a - b)^2\ge 0$, which is always true. Multiply the first inequality by $3$ and we see that it is simply the second inequality with the variables swapped; hence by symmetry it is true as well.

 Equality occurs when $a = b$ and $c = d$, or when $ABCD$ is a kite.

 \textbf{Resources} 


	\item (\textit{USAMO 2004}) Suppose $a_1, \ldots, a_n$ are integers whose greatest common divisor is 1.  Let $S$ be a set of integers with the following properties:


 (a) For $i = 1, \ldots , n$, $a_i \in S$. 

(b) For $i, j = 1, \ldots, n$ (not necessarily distinct), $a_i - a_j \in S$. 

(c) For any integers $x, y \in S$, if $x+y \in S$, then $x-y \in S$.


 Prove that $S$ must be the set of all integers.


 



\textbf{Solution}

Suppose $a_i$ has only one element; then for the greatest common divisor to be 1, $1$ has to be the sole element. Then $1$ is in $S$ by (a), $0$ is in $S$ by (b), $0 + 1 = 1\in S\Rightarrow 0 - 1 = - 1\in S$ by (c), and we can apply (c) analogously to get that $n\cdot 1 \in S$ for integers $n$ and hence $S$ is the set of all integers, as desired.

 \textbf{Lemma:} If $x,y\in a_i$, then $ax + by\in S$ for integers $a,b$.

 \textbf{Proof:} Assume $a_i$ has at least two elements; $x$ and $y$. By (b), $x - y$ is in $S$, and by the application of (c) above, we get that $n(x - y)$ for integers $n$ is in $S$. Then apply (c) to $n(x - y)$ and $ny$ or $nx$ to get that $ax + by\in S$ for all $a,b\in \mathbb{Z}$.

 Now let the terms be $a_1,a_2,\ldots,a_{n}$. By applying our lemma many times, all numbers in  the form $\sum c_1a_1$ for a sequence of integers $c_i$ are attainable if the sequence is of a length which is a power of 2. If not, we "pad" the sequence with many copies of an existing element of the sequence until it does have a length which is a power of 2 - it is apparent that this will not change $S$.

 By Schur's theorem (a generalization of the more well-known Chicken McNugget theorem), every integer greater than some integer $n$ is attainable, and hence there are two members of $S$ in the form $\sum c_1a_1$ which are consecutive integers. Furthermore, because such numbers are closed under addition, their sum is in $S$, and hence so is their difference; $1$. Thus, by the argument at the beginning at this proof, $S$ is the set of all integers, as desired.

 \textbf{Resources} 


	\item (\textit{USAMO 2004}) For what values of $k > 0$ is it possible to dissect a $1 \times k$ rectangle into two similar, but incongruent, polygons?


 



\textbf{Solution}

We will show that a dissection satisfying the requirements of the problem is possible if and only if $k\neq 1$.

 We first show by contradiction that such a dissection is not possible when $k = 1$. Assume that we have such a dissection. The common boundary of the two dissecting polygons must be a single broken line connecting two points on the boundary of the square (otherwise either the square is subdivided in more than two pieces or one of the polygons is inside the other). The two dissecting polygons must have the same number of vertices. They share all the vertices on the common boundary, so they have to use the same number of corners of the square as their own vertices. Therefore, the common boundary must connect two opposite sides of the square (otherwise one of the polygons will contain at least three corners of the square, while the other at most two). However, this means that each of the dissecting polygons must use an entire side of the square as one of its sides, and thus each polygon has a side of length 1. A side of longest length in one of the polygons is either a side on the common boundary or, if all those sides have length less than 1, it is a side of the square. But this is also true of the polygon, which means that the longest side length in the two polygons is the same. This is impossible since they are similar but not congruent, so we have a contradiction.

 We now construct a dissection satisfying the requirements of the problem when $k\neq 1$. Notice that we may assume that $k > 1$, because a $1\times k$ rectangle is similar to a $1\times\frac{1}{k}$ rectangle.

 We first construct a dissection of an appropriate chosen rectangle (denoted by $ABCD$ below) into two similar noncongruent polygons. The construct depends on two parameters ($n$ and $r$ below). By appropriate choice of these parameters we show that the constructed rectangle can be made similar to a $1\times k$ rectangle, for any $k > 1$. The construction follows.

 Let $r > 1$ be a real number. For any positive integer $n$, consider the following sequence of $2n + 2$ points:
 \[A_0 = (0,0), A_1 = (1,0), A_2 = (1,r), A_3 = (1 + r^2, r), \\ A_4 = (1 + r^2, r + r^3), A_5 = (1 + r^2 + r^4, r + r^3),\]and so on, until
 \[A_{2n+1} = (1 + r^2 + r^4 + \cdots + r^{2n}, r + r^3 + r^5 + \cdots + r^{2n - 1}).\]Define a rectangle $ABCD$ by
 \[A = A_0, B = (1 + r^2 + \cdots + r^{2n}, 0), C = A_{2n + 1}, \text{ and }D = (0, r + r^3 + \cdots + r^{2n - 1}).\]The sides of the $(2n+2)$-gon $A_1A_2\ldots A_{2n+1}B$ have lengths
 \[r, r^2, r^3, \ldots, r^{2n}, r + r^3 + r^5 + \cdots + r^{2n-1}, r^2 + r^4 + r^6 + \cdots + r^{2n},\]and the sides of the $(2n+2)$-gon $A_0A_1A_2\ldots A_{2n}D$ have lengths
 \[1, r, r^2, \ldots, r^{2n-1}, 1 + r^2 + r^4 + \cdots + r^{2n-2}, r + r^3 + r^5 + \cdots + r^{2n-1},\]respectively. These two polygons dissect the rectangle $ABCD$ and, apart from orientation, it is clear that they are similar but noncongruent, with coefficient of similarity $r > 1$. The rectangle $ABCD$ and its dissection are thus constructed.

 The rectangle $ABCD$ is similar to a rectangle of size $1\times f_n(r)$, where
 \[f_n(r) = \frac{1 + r^2 + \cdots + r^{2n}}{r + r^3 + \cdots + r^{2n-1}}.\]It remains to show that $f_n(r)$ can have any value $k > 1$ for appropriate choices of $n$ and $r$. Choose $n$ sufficiently large so that $1 + \frac{1}{n} < k$. Since
 \[f_n(1) = 1 + \frac{1}{n} < k < k\frac{1 + k^2 + \cdots + k^{2n}}{k^2 + k^4 + \cdots + k^{2n}} = f_n(k)\]and $f_n(r)$ is a continuous function for positive $r$, there exists an $r$ such that $1 < r < k$ and $f_n(r) = k$, so we are done.

 <i>Alternate solutions are always welcome. If you have a different, elegant solution to this problem, please \href{https://artofproblemsolving.com/wiki/index.php?title=2004\_USAMO\_Problems/Problem\_3&action=edit}{add it} to this page.</i>

 \textbf{Resources} 


	\item (\textit{USAMO 2004}) Alice and Bob play a game on a 6 by 6 grid.  On his or her turn, a player chooses a rational number not yet appearing on the grid and writes it in an empty square of the grid.  Alice goes first and then the players alternate.  When all squares have numbers written in them, in each row, the square with the greatest number is colored black.  Alice wins if she can then draw a line from the top of the grid to the bottom of the grid that stays in black squares, and Bob wins if she can't.  (If two squares share a vertex, Alice can draw a line from on to the other that stays in those two squares.  Find, with proof, a winning strategy for one of the players.


 



\textbf{Solution 1}

Before the game, Bob may select three \textit{useless} squares per row.  He may then move according to the following rules:

 
                \begin{itemize}
                    \item If Alice writes a number in a useless square, then Bob writes a higher number in a non-useless square in the same row on his next turn.\item If Alice writes a number in a non-useless square, then Bob writes a lower number in a useless square in the same row.
                \end{itemize}


                Bob can always do this because there are an even number of squares in each row.

 In this way, Bob can prevent any useless square from being colored black at the end.  Then it will be sufficient for Bob to make the squares (1,1), (2,1), (2,2), (3,2), (3,3), (4,3), (4,4), (5,4), (5,5), (6,5), (6,6) useless, where $(m,n)$ is the square in the $m$th column and $n$th row.  Thus Bob can always win.

 Note that it is not necessary for Bob to move first; even if he moves second, he will always be able to ensure that at least two squares per row of his choosing are not colored black.

 



\textbf{Solution 2}

Bob may move according to the following rules:

 
                \begin{itemize}
                    \item If Alice writes a number in rows 3, 4, 5, or 6, then Bob also writes a number in row 3, 4, 5, or 6.\item If Alice writes a number in row $a \in \{ 1, 2 \}$, column $c$, then Bob writes a number in row $3 - a$, column $c + 3$ (mod 6).  If Alice's number is the highest in its row, then so is Bob's; if it is not, then neither is Bob's.
                \end{itemize}


                In this way, there will be 2 columns in between the black squares in rows 1 and 2, so they cannot possibly touch, which means that Bob will win.

 
<i>Alternate solutions are always welcome. If you have a different, elegant solution to this problem, please \href{https://artofproblemsolving.com/wiki/index.php?title=2004\_USAMO\_Problems/Problem\_4&action=edit}{add it} to this page.</i>

 \textbf{Resources} 


	\item (\textit{USAMO 2004}) Let $a$, $b$, and $c$ be positive real numbers.  Prove that


 
$(a^5 - a^2 + 3)(b^5 - b^2 + 3)(c^5 - c^2 + 3) \ge (a+b+c)^3$.


 


 


	\item (\textit{USAMO 2004}) A circle $\omega$ is inscribed in a quadrilateral $ABCD$.  Let $I$ be the center of $\omega$.  Suppose that


 
$(AI + DI)^2 + (BI + CI)^2 = (AB + CD)^2$.


 


 Prove that $ABCD$ is an \href{https://artofproblemsolving.com/wiki/index.php?title=Isosceles\_trapezoid}{isosceles trapezoid}.


 



\textbf{Solution}

Our proof is based on the following key Lemma.

 \textbf{Lemma}: If a circle $\omega$, centered at $I$, is inscribed in a quadrilateral $ABCD$, then
 \[BI^2 + \frac{AI}{DI}\cdot BI\cdot CI = AB\cdot BC.\qquad\qquad (1)\]

 
                    \begin{figure}[H]
                        \centering
                        \includegraphics[width=0.5\textwidth]{Solution/USAMO_I_2004_6_Solution_Figure_1.png}
                    \end{figure}
                    

 \textit{Proof}: Since circle $\omega$ is inscribed in $ABCD$, we get $m\angle DAI = m\angle IAB = a$, $m\angle ABI = m\angle IBC = b$, $m\angle BCI = m\angle ICD = c$, $m\angle CDI = m\angle IDA = d$, and $a + b + c + d = 180^\circ$. Construct a point $P$ outside of the quadrilateral such that $\triangle ABP$ is similar to $\triangle DCI$. We obtain
 \begin{align*} m\angle PAI + m\angle PBI &= m\angle PAB + m\angle BAI + m\angle PBA + m\angle ABI \\ &= m\angle IDC + a + m\angle ICD + b \\ &= a + b + c + d = 180^\circ, \end{align*}
implying that the quadrilateral $PAIB$ is cyclic. By Ptolemy's Theorem, we have $AI\cdot BP + BI\cdot AP = AB\cdot IP$, or
 \[BP\cdot\frac{AI}{IP} + BI\cdot\frac{AP}{IP} = AB.\qquad\qquad (2)\]Because $PAIB$ is cyclic, it is not difficult to see that, as indicated in the figure, $m\angle IPB = m\angle IAB = a$, $m\angle API = m\angle ABI = b$, $m\angle AIP = m\angle ABP = c$, and $m\angle PIB = m\angle PAB = d$. Note that $\triangle AIP$ and $\triangle ICB$ are similar, implying that
 \[\frac{AI}{IP} = \frac{IC}{CB}\text{  and  }\frac{AP}{IP} = \frac{IB}{CB}.\]Substituting the above equalities into the identity $(2)$, we arrive at
 \[BP\cdot\frac{CI}{BC} + \frac{BI^2}{BC} = AB,\]or
 \[BP\cdot CI + BI^2 = AB\cdot BC.\qquad\qquad (3)\]Note also that $\triangle BIP$ and $\triangle IDA$ are similar, implying that $\frac{BP}{BI} = \frac{IA}{ID}$, or
 \[BP = \frac{AI}{ID}\cdot IB.\]Substituting the above identity back into $(3)$ gives the desired relation $(1)$, establishing the Lemma. $\blacksquare$

 Now we prove our main result. By the Lemma and symmetry, we have
 \[CI^2 + \frac{DI}{AI}\cdot BI\cdot CI = CD\cdot BC.\qquad\qquad (4)\]Adding the two identities $(1)$ and $(4)$ gives
 \[BI^2 + CI^2 + \left(\frac{AI}{DI} + \frac{DI}{AI}\right)BI\cdot CI = BC(AB + CD).\]By the AM-GM Inequality, we have $\frac{AI}{DI} + \frac{DI}{AI}\geq 2$. Thus,
 \[BC(AB + CD)\geq IB^2 + IC^2 + 2IB\cdot IC = (BI + CI)^2,\]where the equality holds if and only if $AI = DI$. Likewise, we have
 \[AD(AB + CD)\geq (AI + DI)^2,\]where the equality holds if and only if $BI = CI$. Adding the last two identities gives
 \[(AI + DI)^2 + (BI + CI)^2\leq (AD + BC)(AB + CD) = (AB + CD)^2,\]because $AD + BC = AB + CD$. (The latter equality is true because the circle $\omega$ is inscribed in the quadrilateral $ABCD$.)

 By the given condition in the problem, all the equalities in the above discussion must hold, that is, $AI = DI$ and $BI = CI$. Consequently, we have $a = d$, $b = c$, and so $\angle DAB + \angle ABC = 2a + 2b = 180^\circ$, implying that $AD\parallel BC$. It is not difficult to see that $\triangle AIB$ and $\triangle DIC$ are congruent, implying that $AB = CD$. Thus, $ABCD$ is an isosceles trapezoid.

 <i>Alternate solutions are always welcome. If you have a different, elegant solution to this problem, please \href{https://artofproblemsolving.com/wiki/index.php?title=2004\_USAMO\_Problems/Problem\_6&action=edit}{add it} to this page.</i>

 \textbf{Resources} 


\end{enumerate}